% Source: https://artofproblemsolving.com/community/c3h1374100_preparation_contest
% Stable identity: aops-topic-1374100-q01
A radius that bisects a chord (not a diameter) is perpendicular to it, so $OB\perp CH$ and $CB=\dfrac{6}{2}=3$.

\begin{asy}
usepackage("amsmath");
size(4.8cm);
defaultpen(fontsize(10pt));
pair ft(pair P,pair U,pair V){return U+dot(P-U,V-U)/dot(V-U,V-U)*(V-U);}
path ra(pair P,pair u,pair v,real s=0.2){return P+s*unit(u)--P+s*unit(u)+s*unit(v)--P+s*unit(v);}
pair O=(0,0),A=(0,5),B=(0,4),C=(-3,4),H=(3,4);
draw(circle(O,5));
draw(C--H,linewidth(0.9)); draw(O--A,linewidth(0.9)); draw(O--C,gray+dashed);
draw(ra(B,H-B,O-B,0.4));
dot(O);dot(A);dot(B);dot(C);dot(H);
label("$O$",O,S); label("$A$",A,N); label("$B$",B,NE); label("$C$",C,W); label("$H$",H,E);
label("$4$",(0,2),E); label("$3$",(-1.5,4),S); label("$5$",(O+C)/2,SW,gray);
label("\mbox{}",(min(currentpicture,user=true)+max(currentpicture,user=true))/2,black); // keeps the text after the figure black when LaTeX typesets the labels inline
\end{asy}

In right triangle $OBC$, the radius is $OC=\sqrt{4^2+3^2}=5$. So $OA=5$ and
\[
BA=OA-OB=5-4=\boxed{1}.
\]
